% Transcription — fonds Alexandre Grothendieck, shelfmark 45, batch 3 (pages 41-60).
% Established from the facsimile archives/batches/45/batch-03.pdf, cut from a
% copy of 45.pdf fetched through the project's relay
% (https://grothendieck-relay.onrender.com/source/45.pdf); PDF page k =
% archivists' page 40+k, no cover sheet. Per the skill transcribe-grothendieck.
% The folder is seventy-nine pages in four batches (1-20, 21-40, 41-60,
% 61-79); this is the third, the others are done by separate passes.
%
% Pass: Opus 5.5 (claude-opus-5-5), 2026-09-26 — first pass, unchecked against the pages by a human.
%
% Covers, no \page marker: page 43, a tan sheet blank but for « H^1 » top
% right; page 60, a tan sheet blank but for three lines top right, read
% « Formes quadratiques / bilinéaires symétriques / etc » (first two words
% doubtful). Both in his hand; their words become section headings, with a
% \note each. Page 60 closes the batch: the run it names begins in batch 4.
%
% Pages skipped, without description: 42, 45, 47, 49, 51, 53, 55, 57, 59.
% They are typed versos (reused paper) carrying nothing in his hand: 42 is
% page 23, « BIBLIOGRAPHIE », of a typescript citing Arf, Artin, Brauer,
% Chevalley-Weil, Deuring, Serre (Corps locaux, 1959), Swan (Induced
% representations, « à paraître »), Witt — on ramification and group rings,
% not semisimple groups; 45 and 47 (scanned upside down) are stencilled
% pages 73 and 83 of a text numbered « n° 239 », § 4 « Prolongement de
% sections. Faisceaux fins »; 49, 51, 53, 55, 57 and 59 (the first four upside
% down) are stencilled pages 59-60, 68 and 73-75 of a text numbered
% « n° 269 » on hermitian, sesquilinear and quadratic forms (isotropic
% subspaces, « Formes bilinéaires associées à une forme hermitienne »).
% Their underlinings are stencilled (the violet of the text), not his; the
% only other marks are the archivists' pencilled page numbers and, on 45 and
% 59, a short stray ink stroke in the left margin, which is not a word or a
% sign and is taken for show-through of his notes on the facing recto.
% Pages summarised: none; page 58, a column of statements linked by double
% arrows and grouped by braces, is given top to bottom with a \note saying so,
% and page 52's scheme of implications is given in two lines with a \note on
% the arrows around it.
%
% The 1961 annotated typescript of the folder title does not appear in pages
% 41-60 either. The nine typed pages here are not it: none carries an
% annotation (see above); they are scrap paper under his own notes, each the
% verso of a leaf of his notes (the typed text shows through several of the
% rectos, e.g. 46, 48, 56); and none is about semisimple groups (sheaves,
% hermitian forms, a bibliography on ramification). So no decision on its
% authorship is made in this batch; it should be looked for in batch 4
% (61-79).
%
% Pages transcribed: 41, 44, 46, 48, 50, 52, 54, 56, 58. His pagination:
% none; archivists' numbers only. Page 41 continues the Lie-algebra
% questions of batch 2 (normalisers in the Lie algebra of Sl(2) in
% characteristic 2); from page 44 on, the pages run as one continuous text on
% forms of a reductive group, page 50 taking up page 48's last sentence.
%
% Notation, kept from batches 1-2: gothic g, h, L as \mathfrak{g},
% \mathfrak{h}, \mathfrak{L}; \underline{\mathrm{Aut}}, Norm. New in this
% batch: G_0 the reductive group « type » over Z, Gamma_0 =
% \underline{Aut}(G_0), Pi_0 = ad(G_0), E_0 = \underline{Ext}(G_0) (Gamma_0
% = E_0 . Pi_0), T_0 ⊂ B_0; rho the element of H^1(S, E_0), G_1 = G_0^rho
% (the « forme ... type », middle word unread), Pi_1, Gamma_1, E_1 their
% twists, T_1 ⊂ B_1 ⊂ G_1, N_1 the normaliser of T_1, W_1 (= W_0) the Weyl
% group; xi ∈ H^1(S, ·) the class of G, \bar\xi its image; P/N_1, P/B_1,
% P/T_1 the schemes C(G), D(G), C'(G) (\mathcal{C}, \mathcal{D},
% \mathcal{C}') of maximal tori, Borels and Killing couples.
%
% What the batch is about. Page 41: Lie algebras of Sl(2) and of Spl(2) in
% characteristic 2, and the normalisers of k.Y, k.H in them. Pages 44-58:
% the classification of forms of a reductive group scheme G_0 over a base S.
% A form G corresponds to a torsor under Gamma = Aut(G_0); its image rho in
% H^1(S, E_0) fixes the « outer type », and the forms with given rho are the
% twists of the quasi-split model G_1 = G_0^rho by H^1(S, Pi_1) (« formes
% intérieures relatives à rho »), with the exact sequences in H^0, H^1 of
% 1 -> Pi -> Gamma -> E -> 1 and the question whether Aut(G) -> Ext(G) is
% onto (false over R, hermitian forms). Pages 48-52: the fibres of
% H^1(S, B_1), H^1(S, N_1), H^1(S, \bar N_1), H^1(S, T_1) -> H^1(S, G_1) as
% Borels, maximal tori, Killing couples (T ⊂ B) up to conjugacy; « tores de
% Killing », interior and extensible isomorphisms T_1 -> T; a Théorème with
% equivalent conditions (i)-(iii') for G to admit a Killing couple, and a
% N.B. on the split case, Isom_S(T_1, T) finite étale under W, and
% Pic(S) = e. Page 54: Questions 1-3 (is every torus contained in a Borel
% interior? does « every maximal torus is in a Borel » force G ≃ G_1? true
% over a field with H^1(k, T) = 0; doubt over R, unitary groups). Page 56:
% conjugacy of Killing couples as injectivity of H^1(S, T_1) -> H^1(S, G_1),
% with GL(2), E = W_0 = Z/2Z as test case. Page 58: the table of
% equivalences — H^1(S, G_1) -> H^1(S, Gamma_1) injective iff Aut(G) ->
% Ext(G) onto; H^1(S, B_1) -> H^1(S, G_1) injective iff any two Borels are
% conjugate; the same for Cartan subgroups and H^1(S, N(T)) — each with his
% verdict on its truth.

\documentclass[11pt,a4paper]{article}
% Le préambule commun à toutes les transcriptions du fonds Grothendieck.
%
% Il tient en une idée : **l'incertitude se compose, elle ne se résout pas.**
% Une transcription qui lisse un mot douteux a détruit la seule chose pour
% laquelle on la faisait — un lecteur doit pouvoir savoir, à chaque endroit,
% ce qui était sur le feuillet et ce qui a été ajouté. Les six macros qui
% suivent sont donc l'appareil critique, et non de la décoration.
%
% Les mêmes commandes sont reconnues par `scripts/render.mjs`, qui produit la
% vue de lecture du site. Une macro ajoutée ici doit l'être aussi là-bas,
% faute de quoi le rendu échouera bruyamment — ce qui est voulu : un rendu
% silencieusement faux est pire qu'un rendu absent.

\ProvidesPackage{grothendieck}[2026/08/08 Transcriptions du fonds Grothendieck]

\usepackage{amsmath,amssymb,amsthm}
\usepackage{mathrsfs}
\usepackage{stmaryrd}
% Les diagrammes commutatifs sont partout dans ce fonds : c'est la langue
% même de la géométrie algébrique de Grothendieck, et une transcription qui
% les rendrait en prose ne transcrirait rien.
\usepackage{tikz-cd}
% Français par défaut — c'est la langue des feuillets — mais l'anglais est
% chargé aussi : la traduction et le résumé sont des documents anglais, et la
% typographie française (espaces avant les deux-points) y serait une faute.
% Ces documents basculent par \selectlanguage{english} en tête de corps.
\usepackage[english,french]{babel}
\usepackage{fontspec}
\usepackage{geometry}
\geometry{a4paper,margin=25mm,marginparwidth=28mm}
\usepackage{marginnote}
\usepackage{xcolor}
% ulem, not soul. `soul`'s \st and \ul are built on a character-by-character
% scan that XeLaTeX's font machinery breaks: the document fails with an
% undefined control sequence rather than degrading. ulem does the same two jobs
% and is Unicode-engine safe. `normalem` stops it hijacking \emph, which the
% transcriptions use for Grothendieck's own emphasis.
\usepackage[normalem]{ulem}
\usepackage{hyperref}
\hypersetup{colorlinks=true,linkcolor=black,urlcolor=black,pdfborder={0 0 0}}

% --- Métadonnées du lot -----------------------------------------------------
% Reprises telles quelles par le script de rendu, qui les affiche en tête de
% la vue de lecture. Elles fixent ce que le fichier prétend couvrir : sans
% elles, une transcription flotte sans point d'attache dans un fonds de seize
% mille feuillets.

\newcommand{\folder}[1]{\gdef\grothFolder{#1}}
\newcommand{\batch}[1]{\gdef\grothBatch{#1}}
\newcommand{\leaves}[2]{\gdef\grothFirst{#1}\gdef\grothLast{#2}}
\newcommand{\foldertitle}[1]{\gdef\grothFoldertitle{#1}}
\gdef\grothFolder{?}\gdef\grothBatch{?}\gdef\grothFirst{?}\gdef\grothLast{?}\gdef\grothFoldertitle{}

% --- Le résumé --------------------------------------------------------------
%
% Il ouvre la lecture modernisée, sous le titre « Résumé », et il est composé
% en petit. Ce n'est pas un ornement : c'est le seul passage du document qui
% ne soit pas des mathématiques, et un lecteur doit pouvoir le reconnaître
% avant de l'avoir lu — puis le sauter s'il connaît déjà le sujet, ou s'y
% arrêter s'il ne le connaît pas. Le filet qui le suit marque la reprise du
% corps.

\newenvironment{resume}
  {\small\setlength{\parskip}{0.45em}}
  {\par\medskip\hrule\medskip}

% --- Mots-clés ---------------------------------------------------------------
%
% En anglais, en clôture du résumé de la lecture modernisée : le vocabulaire
% moderne sous lequel chercher ce que les feuillets construisent. Le manifeste
% du site (`npm run manifest`) les extrait pour étiqueter la cote entière —
% cette ligne est la source unique des tags, il n'y a pas de fichier à côté.
\newcommand{\keywords}[1]{\par\smallskip\noindent
  {\footnotesize\textsc{Keywords} --- \textit{#1}}\par}

% --- L'appareil critique ----------------------------------------------------

% Le numéro de feuillet, en marge. C'est l'ancre qui permet de régler un
% désaccord sur une formule en regardant *un* feuillet plutôt qu'une chemise
% de six cents. Le site s'en sert aussi pour tourner les pages du fac-similé
% quand on fait défiler la transcription.
\newcommand{\leaf}[1]{%
  \marginnote{\footnotesize\color{gray!70}\textsf{#1}}%
  \label{leaf:#1}\ignorespaces}

% Illisible : on ne devine pas. Un mot inventé qui se lit comme les autres est
% le pire résultat possible.
\newcommand{\ill}{\textcolor{red!60!black}{[\,\ldots\,]}}

% Lecture douteuse : le mot est donné, et signalé comme incertain.
\newcommand{\uncertain}[1]{\uline{#1}}

% Ajout de l'éditeur — une lettre manquante, un mot sous-entendu.
\newcommand{\add}[1]{\textcolor{blue!50!black}{[#1]}}

% Biffé par Grothendieck lui-même. Ce qu'il a rayé fait partie du document :
% une rature dit souvent où la pensée a changé de direction.
\newcommand{\struck}[1]{\sout{#1}}

% Note du transcripteur, sur l'état matériel du feuillet.
\newcommand{\note}[1]{\par\smallskip{\small\color{gray!60!black}\textit{[#1]}}\par\smallskip}

% Note marginale de Grothendieck, distincte du corps du texte.
\newcommand{\marginal}[1]{\par\smallskip\noindent\hspace{1em}%
  {\small\color{gray!60!black}#1}\par\smallskip}

% --- Titre ------------------------------------------------------------------

\AtBeginDocument{%
  \begin{center}
    {\large\bfseries Fonds Alexandre Grothendieck --- cote n\textsuperscript{o}~\grothFolder}\\[2pt]
    {\itshape\grothFoldertitle}\\[4pt]
    {\small feuillets \grothFirst--\grothLast\ (lot \grothBatch)}\\[2pt]
    {\footnotesize\color{gray!60!black}Transcription établie d'après le fac-similé numérisé
     par l'Université de Montpellier.\\
     Les lectures incertaines sont soulignées, les illisibles notées [\,\ldots\,],
     les ajouts entre crochets.}
  \end{center}
  \vspace{1em}}

\endinput


\folder{45}
\batch{3}
\pages{41}{60}
\dating{1961-[vers 1970]}
\watermark{Édition de démonstration}
\foldertitle{Groupes semi-simples : notes manuscrites (s.d.), tapuscrit annoté (1961)}

\begin{document}

\page{41}

\[
X \quad Y \quad H
\qquad
\left\lbrace
\begin{array}{l}
[X, H] = \lambda(H) X \\
{[Y, H]} = -\lambda(H) Y \\
{[X, Y]} = H
\end{array}
\right.
\qquad
T \longrightarrow T'
\]
\note{en haut à droite, sous une flèche verticale, « $T \to T'$ » suivi d'un
signe noirci suivi d'un indice $2$, non lu ; le $H$ de $[X,Y] = H$ surcharge
une autre lettre}

\struck{$[X, H]$} \quad $\mathcal{U} \subset$

\emph{$Sl(2)$}
\note{souligné deux fois ; la lecture $Sl(2)$ est douteuse}

\uncertain{Groupe} \struck{\ill{}} \ill{} \ill{} = \ill{} \ill{} \ill{} \ill{}.

$X \quad H \quad Y$ \quad [$X$ et $H$ réunis par une accolade]
\qquad $[X, Y] = 0$

Algèbre de Lie de $Sl(2, k)$ \quad \uncertain{car.~$2$}.
\struck{$[X, H] = \lambda \ldots$}
\[
\left\lbrace
\begin{array}{l}
[H, X] = 0 \\
{[H, Y]} = 0 \\
{[X, Y]} = H
\end{array}
\right.
\qquad \mathcal{U}
\]
\note{les seconds membres des deux premières lignes sont d'abord écrits puis
noircis, et remplacés par $0$ ; « car. 2 » est une lecture douteuse}

Normalisateur de $k.Y = \mathcal{U}$ \ldots\ $k.Y + k.H = \mathfrak{L}$

Normalisateur de $k.H = \mathfrak{h}$ \ldots\ $\mathfrak{g}$ \quad
\uncertain{tt} \ill{}
\note{le $Y$ de $k.Y$ surcharge un $H$ noirci ; les deux « \ldots » tiennent
lieu d'une courte flèche}

Algèbre de Lie de $Spl(2, k)$
\note{« $Spl$ » surcharge un autre mot ; lecture douteuse}
\[
\left\lbrace
\begin{array}{l}
[H', X] = X \\
{[H', Y]} = -Y \\
{[X, Y]} = 0
\end{array}
\right.
\]
Normalisateur de \ill{} $k.Y$ \ldots\ $\mathfrak{g}$

Normalisateur de \ill{}
\note{la page s'arrête sur cette ligne inachevée}


\subsection*{$H^1$}

\note{inscrit en haut à droite d'un feuillet par ailleurs blanc (page 43),
de sa main ; le feuillet ne reçoit pas de numéro de page}

\page{44}

\subsection*{Forme \ill{} type}
\note{titre souligné ; le mot central, qui revient trois fois plus bas
(« forme \ill{} type »), n'est pas lu}

$G_0$ un \uncertain{groupe} \uncertain{réductif} type (\uncertain{sur}
$\mathbb{Z}$), $\Gamma_0 = \underline{\mathrm{Aut}}_{\mathbb{Z}}(G_0)$,
\struck{\ill{}}, \ldots\ $E_0$ \uncertain{le groupe} \ill{}, $T_0 \subset B_0$,
$S$ \uncertain{schéma} de base.
\marginal{$E_0 = \underline{\mathrm{Ext}}(G_0)$, \quad
$\Pi_0 = \mathrm{ad}(G_0)$ \qquad $\Gamma_0 \simeq E_0 \cdot \Pi_0$}
\note{en haut à droite, dans un cadre ; la dernière formule est écrite en
travers d'un passage noirci}

$G$ forme de $G_0$ sur $S$, \uncertain{correspondant} à un
\uncertain{principal} \uncertain{homogène} \ill{} sous $\Gamma_S$,
\struck{\ill{}} le principal homogène \uncertain{associé} sous $E$, qui
\ill{} \uncertain{par} \ill{} de $\rho$ \ill{} \uncertain{l'élément}
\uncertain{correspondant} de $H^1(S, E_0)$
[\uncertain{par} \ill{} \ill{}, \ill{} $\rho = \ill{}$],
\struck{\ill{}}
\note{au-dessus de « le principal », une surcharge noircie portant $P_0$
et $P$ ; un « $P_0$ » et un « $\Gamma_S$ » isolés au-dessus de la ligne}

$G_1 = G_0^{\rho}$ déduit de $G_0$ \ill{} \uncertain{dépendant}
\uncertain{de} $\rho$ : une telle \struck{\ill{}} forme \ill{} \uncertain{dite}
\emph{forme \ill{} type} associée à $\rho$.

[\uncertain{On dit} \uncertain{que} $G$ est une forme \ill{} type s'il
\ill{} \emph{\uncertain{isom.}}\ à $G_1^{\rho}$, \struck{\ill{}}
\ill{} \ill{} $\rho \simeq \rho_G$]

Si \uncertain{on} \uncertain{part} de $\rho$, [\uncertain{sans} $G$], la
\emph{\uncertain{donnée} d'une $G$ \uncertain{forme} \ill{}}
$\rho \xrightarrow[\sim]{\ \varphi\ } \rho_G$
« \emph{\uncertain{équivalente}} \emph{à la} » \emph{donnée d'un fibré
principal homogène sous} $\Pi_1 = \Pi_0^{\rho}$, \ill{}
\note{les deux « \emph{\ldots} » soulignés le sont d'un trait ondulé ; le
$\Pi_1 = \Pi_0^{\rho}$ surcharge une première écriture}

\uncertain{Classification} \uncertain{des} \uncertain{objets} $G$ à
\uncertain{isom.}\ \uncertain{près}
[\ill{} \ill{} \ill{} $(G, \varphi)$ et $(G', \varphi')$ \uncertain{étant}
\ill{} \ill{} \uncertain{un} \uncertain{isom.}\ $\psi : G \to G'$ tel que
l'\uncertain{isom.}\ $\rho_G \xrightarrow{\sim} \rho_{G'}$ \ill{}
\uncertain{compatible} \ill{} \ill{} $\rho_G$ et $\rho_{G'}$ \ill{} $\rho$ ;
\emph{\ill{} \ill{}} \ill{} \ill{} \uncertain{dits} \emph{intérieurs}]
\struck{\uncertain{vérifiant}} \struck{\ill{} \ill{}}
\ill{} \uncertain{donnés} \uncertain{par}
$H^1(S, \Pi_1)$. \uncertain{On} \ill{} \ill{}
\[
H^1(S, \Pi_1) \longrightarrow H^1(S, \Gamma)
= \text{\uncertain{ens.}\ des classes d'\uncertain{isom.}}
\]
\uncertain{vues} \uncertain{des} formes de $G$ \uncertain{sur} $S$ ;
\uncertain{son} \uncertain{ensemble} \ill{} : \uncertain{isom.}\ \uncertain{près}
\uncertain{de} « formes intérieures relatives à $\rho$ ».
\note{le $\Pi_1$ de ces deux $H^1$ surcharge un premier symbole noirci ; la
fin de la phrase, écrite sous la formule, est d'une lecture très incertaine}

\marginal{\uncertain{Distinguer} \ill{} \ill{} \ill{} \ill{} \ill{}
$G_0$ et $\Pi_0$, $G$ et $\Pi_1$ \ldots\
\uncertain{Réalisation} \ill{} \uncertain{la} \uncertain{notion}
\uncertain{d'isom.}\ \emph{intérieur} \ill{}}
\note{deux notes écrites en biais dans la marge gauche ; lecture très
incomplète}

\page{46}

\uncertain{La} \uncertain{fibration} \ill{} \ill{} \ill{} \ill{} $\rho$
\ill{} \ill{} \ill{} \uncertain{groupes} sur $S$
\[
e \longrightarrow \Pi_1^{\rho} \longrightarrow \Gamma_1^{\rho}
\longrightarrow E_1^{\rho} \longrightarrow e
\qquad \Gamma_1^{\rho} \simeq E_1^{\rho} \cdot \Pi_1^{\rho}
\]
\uncertain{qui} \uncertain{fournit} d'\uncertain{ailleurs}
\[
H^0(S, E_1^{\rho}) \longrightarrow H^1(S, \Pi_1^{\rho}) \longrightarrow
H^1(S, \Gamma_1^{\rho}) \longrightarrow H^1(S, E^{\rho})
\]
\[
H^1(S, \Gamma_1^{\rho}) \simeq H^1(S, \Gamma)
\]
\note{les $\Pi$ sont des signes noircis et repassés ; les exposants $\rho$
sont lus sur un signe peu formé}

\struck{On voit} A priori, l'application \uncertain{canonique}
\uncertain{n'est} \uncertain{pas} \uncertain{injective}.
\struck{\ill{}} \uncertain{Deux} \uncertain{éléments}
\struck{\ill{}} \uncertain{de} $H^1(S, \Pi_1^{\rho})$ \uncertain{qui}
\uncertain{donnent} \uncertain{l'élément} \uncertain{neutre} \uncertain{de}
$H^1(S, \Gamma_1^{\rho})$ \struck{\ill{} \ill{} \ill{}} \ill{}
\uncertain{en} \uncertain{correspondance} \uncertain{biunivoque}
\ill{} $H^0(S, E^{\rho}) / H^0(S, \Gamma^{\rho})$,
\uncertain{et} \uncertain{triviale} \uncertain{puisque} $\Gamma_1^{\rho}$
\ill{} \ill{} \uncertain{produit} \uncertain{semi-direct}.
\note{« Deux » est écrit au-dessus d'un mot biffé en fin de ligne}

Mais si $\xi \in H^1(S, \Gamma_1^{\rho})$, \ill{} \uncertain{peut}
\uncertain{utiliser} $\xi$ \uncertain{pour} \uncertain{tordre} \ill{} \ill{}
\uncertain{la} \uncertain{situation}, \ill{} \uncertain{trouve}
\[
e \longrightarrow (\Pi_1^{\rho})^{\xi} \longrightarrow
(\Gamma_1^{\rho})^{\xi} \longrightarrow (E_1^{\rho})^{\xi} \longrightarrow e
\]
\[
e \longrightarrow \mathrm{Ad}(G) \longrightarrow \mathrm{Aut}(G)
\longrightarrow \mathrm{Ext}(G) \longrightarrow e
\]
\note{les deux suites sont reliées par des signes $\simeq$ verticaux ; le
dernier exposant de la première est noirci}

\ill{} \ill{} \ill{} \ill{} \ill{} \uncertain{la} \uncertain{question}
\uncertain{si} \uncertain{tout} \uncertain{automorphisme} \uncertain{ext.}\
de $G$ \uncertain{provient} d'un \uncertain{automorphisme} \ill{}.
\uncertain{C'est} \uncertain{vrai} \uncertain{pour} \uncertain{tout}
\uncertain{\ill{}} \uncertain{évidemment} \uncertain{quand} $G$
\struck{\ill{} \ill{}}
\uncertain{est} \uncertain{une} \uncertain{forme} \ill{} type,
[\uncertain{i.e.}\ \struck{\ill{} \ill{}} \uncertain{possède}
\uncertain{un} \uncertain{sous-groupe} \uncertain{de} Borel $\supset$
\uncertain{Tore} \uncertain{maximal}].

\marginal{\uncertain{exemple} \ill{} \ill{} \ill{} \uncertain{corps}
$\mathbb{R}$ ? \ \ Faux sur $\mathbb{R}$ avec les formes
\uncertain{réelles} (\uncertain{cf.}\ \uncertain{ci-dessous}) : deux formes
hermitiennes \uncertain{donnant} \ill{} \ill{} \ill{} $\mathbb{R}$ \ill{}
mais \uncertain{non} \uncertain{isom.}}
\note{écrit en biais en bas à gauche ; la première ligne est entourée}

\page{48}

\note{la page s'ouvre sur un tableau à trois colonnes, que l'on donne ligne
par ligne}

\uncertain{Schéma} des tores maximaux de $G$ \ldots\
$P/N_1 \simeq Q/N'_1 = \mathcal{C}(G)$

\quad --- \quad Borel \quad --- \quad
$P/B_1 \simeq Q/B'_1 \simeq \mathcal{D}(G)$

\quad --- \quad couples de Killing \quad --- \quad
$P/T_1 \simeq Q/T'_1 \simeq \mathcal{C}'(G)$

\quad --- \quad $\lbrace$tores maximaux \uncertain{munis} de $G$ ; «
\quad » \uncertain{immersions} \emph{intérieures} $T_1 \to G\rbrace$ \quad
$\simeq P/T_1$
\note{le nombre $21$ est écrit à gauche de cette dernière ligne ; les deux
lignes entre accolades sont l'une sous l'autre, reliées par des guillemets
de répétition. Une flèche courbe monte de $\mathcal{C}'(G)$ vers
$\mathcal{C}(G)$, une flèche droite de $\mathcal{C}'(G)$ vers
$\mathcal{D}(G)$}

\marginal{fibré principal sous $W_1$ \quad / \quad fibré \struck{\ill{}}
en \uncertain{fibrés} \uncertain{principaux} \uncertain{sous} $B_1/T_1$
$= \mathcal{C}(B_1)$ \uncertain{associé} \uncertain{au} fibré principal
\uncertain{homogène} \ill{} \uncertain{de} \uncertain{groupe} $B_1$}
\note{à droite du tableau, en face des deux flèches}

\emph{N.B.} Soit $T$ un tore \add{\uncertain{maximal} de $G$} ; il y a
correspondance \uncertain{biunivoque} (\uncertain{fonctorielle en}
$S'/S$) entre les ensembles \uncertain{suivants}
\begin{itemize}
  \item[a)] L'\uncertain{ens.}\ des Borel \uncertain{contenant} $T$
  \item[b)] L'\uncertain{ens.}\ des \uncertain{isom.}\ $T_1 \to T$ qui sont
  \emph{intérieurs}
\end{itemize}
\marginal{$\rbrace\ P/T_1 \simeq P'/T'_1$}
\struck{$P$ \ill{} \uncertain{quelconque} \ill{}}
\uncertain{De} \uncertain{plus}, \struck{\ill{} \ill{} \ill{} \ill{}}
\uncertain{considérons} \uncertain{le} \uncertain{sous-ensemble}
\uncertain{suivant} \uncertain{de} b) :
\struck{\uncertain{il} \uncertain{existe} \ill{} \uncertain{application}}
\begin{itemize}
  \item[b')] L'\uncertain{ens.}\ des \uncertain{isom.}\ $T_1 \to T$ qui sont
  « \uncertain{génériquement} \ill{} »
\end{itemize}
\struck{\ill{} \ill{} \ill{}} \uncertain{Il} \uncertain{existe} \ill{}
\uncertain{projection} de b') sur b).
\marginal{$\rbrace\ P'/T_1$}
\marginal{$W_1$ et $W$ \ill{} \ill{} \ill{} \ill{} \ill{} \ill{} \ill{}
\ill{} \ill{} $\neq \rho$]}
\note{la dernière note est écrite en biais dans la marge gauche}

De plus, \uncertain{les} \uncertain{conditions} \uncertain{suivantes}
\uncertain{sur} $T$ \uncertain{sont} \uncertain{équivalentes}
\begin{itemize}
  \item[$\alpha$)] $T$ \uncertain{est} \uncertain{contenu} \uncertain{dans}
  un Borel
  \item[$\beta$)] Il existe un \uncertain{isom.}\ \emph{intérieur}
  $T_1 \to T$
  \item[$\gamma$)] Il existe un \uncertain{isom.}\ \ill{} $T_1 \to T$
\end{itemize}
\uncertain{On} \uncertain{dit} \uncertain{alors} \uncertain{que} $T$
\uncertain{est} \uncertain{un} « Tore de \uncertain{Killing} ».
\uncertain{Pour} qu'il existe, il \uncertain{faut} et il \uncertain{suffit}
\uncertain{que} $G$ \uncertain{soit} \uncertain{isomorphe} à une forme
\ill{} type.
\note{« Killing » surcharge un autre mot}

L'image inverse de $\xi \in H^1(S, G_1)$ dans $H^1(S, B_1)$ est en
correspondance biunivoque avec l'\uncertain{ens.}\ des sous-groupes de Borel
de $G$, modulo \uncertain{automorphismes} \emph{intérieurs}.
[\uncertain{Je} \uncertain{ne} \uncertain{vois} \uncertain{pas}
\uncertain{d'ailleurs}] [i.e.\ \uncertain{on} \uncertain{peut}
\uncertain{restreindre} \uncertain{le} \uncertain{groupe}
\uncertain{structural} \uncertain{à} $B_1$] \uncertain{ssi}
il existe un \uncertain{sous-groupe} de Borel, \uncertain{dans} \ill{}
\uncertain{cas}, il \ill{} \uncertain{revient} \uncertain{au} \uncertain{même}
\uncertain{de} \uncertain{dire} \uncertain{que} \uncertain{ces} deux
\uncertain{sous-groupes} de Borel \uncertain{de} $G$ \uncertain{sont}
\uncertain{toujours} \uncertain{conjugués}. \struck{\ill{} \ill{}}

\marginal{\ill{} \ill{} \ill{} \uncertain{mais} \ill{} \uncertain{forme}
$G_1$ \uncertain{sur} un type \uncertain{extérieur} \ill{}
\uncertain{mais} \uncertain{seulement} \uncertain{quand} il y a
\uncertain{deux} ($?$) \uncertain{fois} un Borel}
\note{écrit en biais dans la marge gauche}

[\emph{\uncertain{Conjecture} 3} \struck{\ill{}} \uncertain{Contre-exemple}
\ill{} \ill{} \ill{} \ill{} \ill{} \ill{}\,? si $S$ \uncertain{est}
\uncertain{local} ?]
\note{dans un cadre, marqué d'un double trait à gauche et d'une croix à
droite ; un mot entouré, au-dessus, est relié par un trait au
« Contre-exemple »}

\page{50}

L'image inverse de $\xi$ dans $H^1(S, N_1)$ est en correspondance
biunivoque avec l'\uncertain{ens.}\ des Tores maximaux de $G$, modulo
\uncertain{automorphismes} intérieurs. \uncertain{En} \uncertain{particulier},
\struck{\ill{}} \uncertain{i.e.}\ \uncertain{on} \uncertain{peut}
\uncertain{restreindre} \uncertain{le} \uncertain{groupe} \uncertain{structural}
à $N_1$, ssi il existe un \uncertain{tore} \uncertain{maximal} ; $\xi$ s'il
\uncertain{en} \uncertain{est} \uncertain{ainsi}, \uncertain{alors} il \ill{}
\uncertain{identité} \ill{} \ill{} \uncertain{tous} \uncertain{les}
\uncertain{tores} \uncertain{maximaux} \uncertain{sont} \uncertain{conjugués},
[\ill{} \uncertain{qui} \ill{} \ill{} \ill{} \uncertain{exceptionnel}
\ill{} \ill{} \ill{} « \uncertain{arithmétiques} »].
\note{un trait relie « intérieurs » à « restreindre », par-dessus un mot
biffé. Le $\xi$ en début de proposition est surchargé}

\marginal{\uncertain{vaut} \ill{}, \ill{} $W$ $G_1$ \ill{} $W$ \ill{}
\ill{} \ill{} \uncertain{maximal}.}
\marginal{\uncertain{Il} \uncertain{semble} \uncertain{que} \ill{}
\uncertain{soit} \ill{} \ill{} \uncertain{fois} \uncertain{que} \ill{}
$S$ \ill{} \ldots\ ?}
\note{deux notes en marge gauche, la seconde dans un cadre, reliée par une
croix au début de la ligne « tore maximal »}

\struck{\ill{}} Si $S$ local, \struck{\ill{}}, $N_1$ \uncertain{produit}
\uncertain{semi-direct} de $W_1$ \uncertain{et} $T_1$, \uncertain{on a}
$H^1(S, N_1) \twoheadrightarrow H^1(S, W_0)$. \ill{} \ill{} \ill{} \ill{}
\ill{} \ill{} \ill{} \ill{} \ldots\ \uncertain{un} \uncertain{élément} de
$H^1(S, W_0)$ \ill{} \ill{} \ldots
\note{passage ajouté dans l'interligne, en plus petit, pris dans une
accolade à gauche ; très raturé, lecture très incomplète. Au-dessus de
$W_1$ est écrit « $= W_0$ », sous la flèche « sur ». Une bulle à gauche,
reliée à l'accolade, porte « \uncertain{fibrés} \ill{} \ill{} \ill{}
$H^1$ \uncertain{des} $T_1$ \uncertain{tordus} »}

L'image inverse de $\xi$ \uncertain{dans} $H^1(S, \bar N_1)$ est en
correspondance biunivoque avec l'\uncertain{ens.}\ des couples
$(T \subset B)$ de Killing de $G$, \uncertain{modulo} \uncertain{autom.}\
\emph{\uncertain{extérieurs}}. \uncertain{Donc} il \uncertain{est}
\uncertain{non} \uncertain{vide} ssi il existe un tel couple, i.e.\ ssi il
existe un Tore \struck{\uncertain{privilégié}} \uncertain{de} Chevalley,
ou encore un Borel de Chevalley ; [\uncertain{chose}
\uncertain{exceptionnelle} \ill{} \ill{} \uncertain{et} \ill{} \ill{}
\uncertain{arithmétique}] ; \uncertain{et} \struck{\ill{} \ill{}}, s'il en
est ainsi, \uncertain{alors} \uncertain{l'on} \uncertain{a} \ill{}
\uncertain{identité} \uncertain{à} un \uncertain{élément} \uncertain{ssi}
\uncertain{tous} \uncertain{les} couples de Killing de $G$ \uncertain{sont}
\uncertain{conjugués} [\uncertain{et} il \uncertain{est} \uncertain{possible}
\uncertain{que} \uncertain{cette} \uncertain{condition} \uncertain{soit}
\uncertain{vérifiée} \uncertain{si} $S$ \uncertain{est} \uncertain{local}],
\uncertain{i.e.}\ \ill{} \uncertain{le} \uncertain{cas}
\uncertain{où} \ill{} \ill{} \ill{} $T_1$ \uncertain{déployé}, i.e.\ $\rho = e$],
\ill{} \ill{} \ill{} \uncertain{de} \ill{} \ldots]
\note{« Chevalley » est lu deux fois ; le signe de $\rho = e$ est un $e$
entouré}

\marginal{\uncertain{Question} \uncertain{équivalente}
\struck{\ill{}} \uncertain{2} \ill{} \ill{} \uncertain{vérifiée} ; \ill{}
\uncertain{sur} Borel \ill{} \ill{} \uncertain{Dans} \uncertain{tout}
\uncertain{maximal} \uncertain{de} Borel $B$ \ill{} \uncertain{tous}
\uncertain{conjugués} \uncertain{dans} $B$}
\note{en bas à gauche, en biais, dans un cadre ouvert, reliée par une
flèche à la ligne « doit vérifier » ; lecture très incomplète}

\page{52}

\emph{Th.} Conditions équivalentes \uncertain{sur} $\xi \in H^1(S, H_1)$,
\uncertain{définissant} $G = G_1^{\xi}$.
\begin{itemize}
  \item[(i)] $\xi \in \mathrm{Im}\, H^1(S, T_1)$ dans $H^1(S, H_1)$
  \item[(i')] $\bar\xi \in \mathrm{Im}\, H^1(S, T'_1)$ \uncertain{dans}
  $H^1(S, \Gamma_1)$ \ / \uncertain{ou} $\bar\xi_0 \in \mathrm{Im}\,
  H^1(S, T'_0)$ \uncertain{dans} $H^1(S, \Gamma'_0)$
  \item[(i'')] $\bar\xi \in \mathrm{Im}\, H^1(S, T_1)$ dans
  $H^1(S, \Gamma_1)$ \ / \uncertain{ou} $\xi_0$
  \item[(ii)] Il existe un couple de Killing dans $G$, i.e.\ un
  Borel $\supset$ Tore \uncertain{max.}
  \item[(iii)] Il existe un \struck{\ill{} \ill{} \ill{} \ill{}
  \ill{} \ill{}} Tore maximal de $G$, \struck{\uncertain{isom.}\ à $T_1$}
  \uncertain{provenant} d'une \emph{immersion} \emph{extensible}
  $T_1 \to G$
  \item[(iii')] Il existe une \emph{immersion intérieure} de $T_1$ dans
  $G$.
\end{itemize}
\note{les numéros (i), (ii), (iii) et (iii') sont entourés ; (iii) est
en outre marqué d'un crochet et de traits, et une ligne biffée au-dessus
de lui commençait « \uncertain{Tore} » ; « immersion » et « extensible »
sont soulignés, ce dernier deux fois ; « immersion intérieure » est
souligné d'un trait ondulé}

\marginal{De plus, si $T$ est un Tore maximal, conditions équivalentes.
a) $\exists$ Borel $\supset T$ \quad b) $\exists$ \uncertain{isom.}\
$T_1 \to T$ \uncertain{extensible} \quad a) $\Rightarrow$ c) $\exists$
\uncertain{isom.}\ $T_1 \to T$ \uncertain{intérieur}}
\note{écrit à droite, en face de (ii) et (iii) ; « extensible » est au
dessus de « intérieur », comme une variante}

\emph{\uncertain{Démonstration}.}
\[
(iii') \Longleftrightarrow (i) \Longrightarrow (i') \Longleftrightarrow (ii)
\qquad (iii) \Longleftrightarrow (i'')
\]
\note{schéma d'implications sur deux rangées, que l'on résume : en haut
« (iii') $\Leftrightarrow$ (i) $\Rightarrow$ (i') $\Leftrightarrow$ (ii) »,
suivi d'un « $\Leftrightarrow$ (iii') » biffé ; en bas « (iii)
$\Leftrightarrow$ (i'') ». Entre les deux rangées, des flèches obliques
dont plusieurs sont barrées en zigzag, deux portant « triviale » et une
« \uncertain{inutile} » ; une grande courbe « triviale » joint (iii') à
(iii). Des traits obliques renvoient aux lignes de droite :
« \ill{} \uncertain{que} \uncertain{ce} \uncertain{ne} \uncertain{sont}
\uncertain{pas} \uncertain{des} \uncertain{conditions} \ill{} », « \ill{}
\uncertain{des} \ill{} \uncertain{du} \uncertain{nuage}\,? » et, biffé,
« \uncertain{réalisations} de $T_1$ en $\Gamma_1$ \uncertain{et} $T_1$ »}

\emph{N.B.} Supposons que $T_1$ soit $\simeq \mathbb{G}_m^r$
(\uncertain{cas} \uncertain{déployé}). Alors \emph{\uncertain{tout}} tore
$T \subset G$ qui est \struck{\ill{} \ill{} \ill{} \ill{} \ill{}}
\uncertain{contenu} dans un Borel, \struck{\ill{}} i.e.\ il existe une
\uncertain{isom.}\ \emph{intérieure} de $T_1$ avec $T$. \uncertain{Le}
\struck{\ill{}} \uncertain{bien} \uncertain{plus},
\struck{\uncertain{c'est} \uncertain{local}} \uncertain{le}
\uncertain{schéma} \uncertain{des} \uncertain{isom.}\ intérieurs de $T_1$
dans $T$
\uncertain{pour} \uncertain{un} \uncertain{tel}
\uncertain{schéma} \uncertain{de} $\underline{\mathrm{Isom}}_S(T_1, T)
\simeq (\mathbb{Z}^{r^2})_S$, \uncertain{qui} \uncertain{est}
\uncertain{fini} \uncertain{étale} \uncertain{sur} $S$ \uncertain{de}
\uncertain{groupe} $W$, \uncertain{donc} \uncertain{l'est}
\uncertain{aussi} \uncertain{diagonalisable} \ldots.
\emph{\uncertain{Cor}}. Donc, \uncertain{si} \uncertain{on} \uncertain{a}
$\mathrm{Pic}(S) = e$, il existe un \ill{} de \ill{} $G_0$ sur
$G$, \uncertain{transformant} $T_0$ en $T$.
\note{« tout » est souligné ; « \uncertain{cas} \uncertain{déployé} » est
souligné d'un trait ondulé ; « $\mathrm{Pic}(S) = e$ » est précédé d'un
signe « $\Xi$ » ou d'un « Et » surchargé. La dernière ligne, écrite serrée
au bas du feuillet, est d'une lecture incertaine}

\page{54}

\emph{Question 1}. Supposons que \struck{\uncertain{$G$ soit un}} $T$
\struck{\ill{}} $\simeq T_1$, \uncertain{dans} $G$ \uncertain{tore}
\uncertain{maximal} \uncertain{contenu} dans un Borel, i.e.\ \ill{} \ill{}
qu'il existe un \uncertain{isom.}\ \emph{intérieur} $T_1 \to T$ ?? \ J'en
doute.

\emph{Question 2}. Supposons que \struck{\ill{} \uncertain{isom.}\ à $G$
\ill{}} \uncertain{tout} \uncertain{tore} \uncertain{maximal} $T$ \uncertain{soit}
\uncertain{contenu} dans un Borel. S'ensuit-il que $G \simeq G_1$ ?
\struck{\ill{} \ill{}} \struck{\ill{} \ill{}} (\uncertain{p.\ ex.}\ $S$
\uncertain{local}) \struck{\ill{}} \uncertain{C'est} \uncertain{le}
\uncertain{cas} si $\mathrm{Pic}(S) = 0$ et si $\rho = e$ \uncertain{est}
\uncertain{trivial}. Est-ce \ill{} \uncertain{vrai} \uncertain{pour}
un \uncertain{corps} $k$ ($\rho$ \uncertain{non} \uncertain{trivial})
?? \struck{\uncertain{Pour} \uncertain{exemple} $k = \mathbb{R}$,
\uncertain{groupes} \uncertain{unitaires}} (\uncertain{C'est}
\uncertain{évident} \uncertain{si} $k$ \uncertain{est} un
\emph{\uncertain{corps} de \uncertain{Tate}} [$H^1(k, T) = 0$ \uncertain{pour}
\emph{\uncertain{tout}} \uncertain{tore} \uncertain{sur} $k$]. \uncertain{On}
\uncertain{a} \ill{} \uncertain{si} $k = \mathbb{R}$, \struck{\ill{}}
$G_1$ \ill{} \uncertain{groupe} \uncertain{unitaire}
\uncertain{correspondant} à \uncertain{une} \uncertain{forme} \uncertain{non}
\uncertain{anisotrope} : \uncertain{deux} \uncertain{variables} ???
\note{« de \uncertain{Tate} » est souligné ; « \uncertain{unitaire} »,
entouré, est relié par un trait à « \uncertain{correspondant} »}

\marginal{$\rho$ \struck{\ill{}}, \uncertain{on} \uncertain{notera}
\uncertain{que} \uncertain{les} \uncertain{normalisateurs} de
$\Gamma_0$ \uncertain{et} $(B_0, T_0)$ \uncertain{dans} $\Gamma_0$
\struck{\ill{}} $T_0$ \ldots\ \uncertain{mais} \ill{}
\struck{\ill{} \ill{} \ill{} \ill{} \ill{} \ill{} \ill{} \ill{} \ill{}}
\ldots\ \uncertain{Dans} \uncertain{ce} \uncertain{cas} \ill{} « $\xi$ »}
\note{deux notes en biais dans la marge gauche, la première en partie
encadrée et barrée d'un long trait oblique, reliée par un trait à la ligne
« Question 2 » ; lecture très incomplète}

\emph{Question 3}. \uncertain{Supposons} \uncertain{que} \ill{} il existe
un \uncertain{sous-groupe} de Borel $B$ de $G$ [\uncertain{i.e.}\ \uncertain{que}
\uncertain{l'on} \uncertain{peut} \uncertain{restreindre} \uncertain{le}
\uncertain{groupe} \uncertain{structural} de $G_1$ à $B_1$]. \uncertain{Alors}
\uncertain{peut-on} \uncertain{restreindre} \uncertain{le} groupe
structural à \uncertain{celui} de $G_1$ à $T_1$, et \uncertain{a-t-on}
\ill{} \struck{\ill{}} \uncertain{l'on} \uncertain{peut}
\uncertain{trouver} un \ill{} \uncertain{sous-groupe}
\note{la page s'arrête sur cette ligne ; au-dessus de « Question 3 », un
signe noirci}

\page{56}

\marginal{$H^1(S, T_1)$ \qquad $H^1(S, T) \to H^1(S, B)$}
\note{deux formules isolées dans la marge gauche, en face de la première
ligne et du milieu de la page}

\uncertain{classification} \uncertain{des} formes intérieures
\uncertain{munies} d'un couple de Killing.

[\emph{N.B.} Si $G$ est une forme intérieure, \uncertain{deux} couples
$(T, B)$, $(T', B')$ de Killing \uncertain{dans} $G$ \uncertain{définissant}
le \uncertain{même} \uncertain{élément} de $H^1(S, T_1)$ \uncertain{sont}
\uncertain{conjugués} \uncertain{par} \uncertain{un} \uncertain{automorphisme}
\uncertain{intérieur}. \uncertain{Cela} \uncertain{signifie} \uncertain{donc}
\uncertain{que} \uncertain{deux} \uncertain{éléments} \uncertain{distincts}
\uncertain{dans} $H^1(S, \bar B)$ \uncertain{mod}~$B$ \ill{} \ill{}
$T = T'$, \uncertain{c'est-à-dire} \uncertain{le} \uncertain{donné}
\uncertain{de} $\sigma \in W(T)$ \uncertain{qui} \uncertain{transforme} $B$
\uncertain{en} $B'$, si $B'$ \uncertain{est} \uncertain{l'opposé} \uncertain{de}
$B$ \uncertain{rel.}\ à $T$, \uncertain{dans} \uncertain{la} \uncertain{question}
\uncertain{est} \struck{\ill{}} \uncertain{si} \uncertain{un} \ill{}
\uncertain{élément} \uncertain{nul} \uncertain{dans} $H^1(S, T)$.]
\note{« $(T, B)$, $(T', B')$ » est ajouté au-dessus de la ligne ;
« mod $B$ » de même ; un mot biffé et une insertion en chevron sur la ligne
de « question »}

La question des \uncertain{deux} couples de Killing \uncertain{dans} un
\uncertain{même} $G$ \uncertain{sont} \uncertain{toujours} \uncertain{conjugués},
\uncertain{équivaut} \uncertain{donc} à \uncertain{celle} \uncertain{de}
\uncertain{savoir} si $H^1(S, T_1) \to H^1(S, G_1)$ \uncertain{est}
\uncertain{toujours} \uncertain{injective}. \uncertain{Par} un
\uncertain{contre-exemple}, il \uncertain{semble} \uncertain{donc}
\uncertain{que} $H^1(S, T_1) \to H^1(S, N_1)$ \uncertain{n'est}
\uncertain{pas} \uncertain{toujours} \uncertain{injective}. \uncertain{On}
\uncertain{devrait} \uncertain{étudier} \uncertain{si} $H^1(S, T_1) \to{}$
\ill{} \ill{} $=$ \uncertain{image} \uncertain{dans} $H^1(S, N_1)$
\uncertain{des} \uncertain{éléments} \uncertain{conjugués} \uncertain{par}
$H^0(S, W_1)$ \uncertain{opérant} \uncertain{de} \uncertain{façon}
\uncertain{évidente}. \uncertain{Je} \uncertain{prendrai} \uncertain{le}
\uncertain{cas} \uncertain{où} $G_0 = GL(2)$, \uncertain{donc}
$E = \mathbb{Z}/2\mathbb{Z}$, $W_0 = \mathbb{Z}/2\mathbb{Z}$
\uncertain{est} \ill{} \ill{}, \struck{\uncertain{donc} $W_0 \ldots W_0$
\ill{}} $E$ \uncertain{y} \uncertain{opère} \uncertain{trivialement}
(\uncertain{il} \ill{} \uncertain{que} \ill{} !) \uncertain{et}
$W_0 = W_1$, \uncertain{qui} \uncertain{opère} \uncertain{par}
\uncertain{symétrie} \uncertain{sur} $T_1$.
\note{« $GL(2)$ » surcharge une autre écriture ; la page s'arrête ici}

\page{58}

\note{la page est une colonne d'énoncés reliés par des doubles flèches
verticales, groupés par trois accolades à droite, chacune suivie d'un
commentaire ; on la donne de haut en bas}

$H^1(S, G_1) \to H^1(S, \Gamma_1)$ \uncertain{injectif} ?

\struck{$H^1$} \quad $\Updownarrow$

$\underline{\mathrm{Aut}}(G) \to \underline{\mathrm{Ext}}(G)$
\uncertain{admet} \uncertain{toujours} \uncertain{une} \uncertain{section}

$\Updownarrow$

$\mathrm{Aut}(G) \to \mathrm{Ext}(G)$ toujours \uncertain{surjectif}

\marginal{\uncertain{forme} \uncertain{extérieure} $G_1$
\uncertain{donnée}}
\note{à gauche de la première accolade}

\emph{\uncertain{Non}} : \uncertain{faux} \uncertain{en} \uncertain{général}
\uncertain{sur} \uncertain{un} \uncertain{corps}, \uncertain{vrai}
\uncertain{pour} \uncertain{les} \uncertain{formes} \uncertain{intérieures}
($\xi = e$) \uncertain{et} \uncertain{pour} \uncertain{les}
\uncertain{extérieures} \uncertain{de} \uncertain{Tits}\,?
\note{commentaire à droite de la première accolade, précédé d'un grand
signe tracé à la plume ; la première ligne est d'une lecture très
incertaine}

$H^1(S, B_1) \to H^1(S, G_1)$ \uncertain{injectif}

$\Updownarrow$

\uncertain{Deux} sous-groupes de Borel \add{\uncertain{dans} \uncertain{un}
$G$} \uncertain{sont} \uncertain{toujours} conjugués
\note{« dans un $G$ » est entouré au-dessus de la ligne}

$\Updownarrow$

\uncertain{Pour} \uncertain{tt} $G$ et $B$, \struck{\uncertain{le}
\uncertain{groupe} \ill{}} $H^1(S, B) \to H^1(S, G)$ \struck{\uncertain{est}
\uncertain{injectif}}

$\Updownarrow$

\uncertain{Pour} \uncertain{tt} $G$, et $B$, \struck{\ill{} \ill{}}
$H^1(S, B) \to H^1(S, G)$ \uncertain{est} \uncertain{bijectif}
\note{« tt » est souligné deux fois dans la première de ces deux lignes}

\uncertain{vrai} \uncertain{pour} \uncertain{les} \uncertain{formes}
\uncertain{intérieures} ($\xi = e$) \uncertain{et} \uncertain{pour}
\uncertain{les} \ill{} \uncertain{de} \uncertain{Tits}.
\note{commentaire à droite de la deuxième accolade}

\marginal{\uncertain{Si} $\ldots$ \uncertain{de} $\ldots$ \ill{}
\uncertain{extérieur} \ill{} \ill{} \uncertain{si} \ill{} \ill{}
\uncertain{équivalent} \ill{} \ill{} \ill{} $G \to G/B$ \ill{}
\uncertain{triviale}}
\note{longue note en biais dans la marge gauche, en face de la deuxième
accolade ; lecture très incomplète}

$H^1(S, N_1) \to H^1(S, G_1)$ \uncertain{injectif}

$\Updownarrow$

Deux sous-groupes de Cartan \uncertain{dans} \uncertain{un} $G$
\uncertain{toujours} \uncertain{conjugués}

$\Updownarrow$

\uncertain{Pour} \uncertain{tt} $G$ et $T$, le \uncertain{morphisme} de
$H^1(S, N(T)) \to H^1(S, G)$ \uncertain{injectif}

$\Updownarrow$ \quad \ill{} \uncertain{injectif}
\note{la dernière ligne est inachevée : un trait, puis « injectif »}

\uncertain{Semble} \uncertain{faux} \uncertain{même} \uncertain{dans}
\uncertain{le} \uncertain{corps} \uncertain{non} \uncertain{algébrique} :
\uncertain{deux} \uncertain{distinct.}\ \uncertain{possible}, \uncertain{même}
\uncertain{pour} \uncertain{les} \uncertain{formes} \struck{\ill{}}
\uncertain{intérieures} \ldots
\note{commentaire à droite de la troisième accolade}

\subsection*{\uncertain{Formes} \uncertain{quadratiques}, bilinéaires symétriques, etc}

\note{inscrit en haut à droite d'un feuillet par ailleurs blanc (page 60),
de sa main, sur trois lignes ; le feuillet ne reçoit pas de numéro de page.
La suite qu'il annonce commence avec le lot suivant}

\end{document}
